\section{Related work}
\label{sec:related}
\paragraph{Cross-task memory for GUI agents.} Cross-task memories answer how to do a task. AutoDroid~\citep{autodroid}, AppAgent~\citep{appagent}, Agent Workflow Memory~\citep{awm}, Synapse~\citep{synapse}, and the tips and shortcuts of Mobile-Agent-E~\citep{mobileagente} store transition graphs, exemplar trajectories, and workflows that transfer across episodes. Mobile-Agent-E's note-taker is the closest prior design to our need: it records values it judges important as the episode proceeds, a write-ahead note that must guess what the agent will need. Neither variant lets the agent re-read an earlier screen in full.

\paragraph{Within-task compression and lookup.} GUI-Rise~\citep{guirise} maintains a rolling summary and PAL-UI~\citep{palui} lets the agent look back at chosen frames. Our summary condition (C) is the per-screen variant of the former, and our screen index is the gist-and-lookup design that ReadAgent~\citep{readagent} introduced for long documents and that MemGPT~\citep{memgpt} and Generative Agents~\citep{generativeagents} instantiate as paged or scored memory streams. These systems prompt their writers; we train ours against the reader, as document expansion~\citep{doc2query} does, and find no gain (\S\ref{sec:writer}).

\paragraph{Long-context evaluation.} Lost-in-the-middle~\citep{lostmiddle} and RULER~\citep{ruler} show that models use long contexts unevenly and that synthetic needle tasks overstate usable length. Ours is a needle task on real screens with contexts of 1--5K tokens, short enough that full history still helps. Its two pitfalls are specific to mining needles from trajectories, where a needle can recur or sit on the current screen.
